\documentclass[11pt]{article}
\usepackage[margin=1in]{geometry}
\usepackage{amsmath}
\usepackage{graphicx}
\usepackage{tikz-cd}
\usepackage{multicol}

\setlength{\parindent}{0pt}
\setlength{\parskip}{1\baselineskip}

\begin{document}

\section*{Tema: formacion de planetas en nebulosas}

Este documento tecnico presenta una propuesta formal para estudiar la \textbf{formacion de planetas en nebulosas protoplanetarias} usando la base computacional del proyecto BeeFlow.

BeeFlow fue implementado como un motor de simulacion hibrida de agentes + campo + dinamica global, con ejecucion determinista, exportacion de metricas y validaciones de rendimiento. Esta arquitectura se considera apta para reformular el problema planetario como un sistema multi-escala con interacciones locales (microfisica) y variables globales (macroestructura del disco).

\section*{Definicion del problema}

Problema cientifico: explicar como una nebulosa protoplanetaria evoluciona desde polvo y gas hasta estructuras planetesimales y, eventualmente, planetas estables.

Desde el punto de vista de simulacion, se requiere:

\begin{itemize}
  \item representar miles de entidades con estado interno y dinamica local;
  \item modelar campos espaciales continuos (densidad, temperatura, senales);
  \item acoplar escalas temporales rapidas y lentas sin perder estabilidad numerica;
  \item correr experimentos reproducibles con semillas fijas y barrido de parametros.
\end{itemize}

La hipotesis de trabajo del proyecto es que una arquitectura tipo BeeFlow puede cubrir estas necesidades en hardware de consumo (CPU multicore), con tick fijo y presupuesto temporal por paso de simulacion.

\section*{Marco teorico}

El problema de formacion planetaria combina transporte de materia, agregacion, perdida y transiciones de estado. En esta adaptacion se consideran tres niveles teoricos:

\textbf{(1) Nivel local (micro):} interacciones entre unidades discretas (particulas, agregados o ``proto-agentes'').

\textbf{(2) Nivel de campo (meso):} distribucion espacial de magnitudes continuas sobre un grid 2D (por ejemplo, densidad efectiva y gradientes de atraccion).

\textbf{(3) Nivel global (macro):} variables lentas del sistema (reserva total de masa util, forzamiento termico, fase evolutiva del disco).

La dinamica en BeeFlow ya implementa esta separacion por capas. Su ecuacion local de actualizacion (en el proyecto original) tiene forma:

\[
E_{t+1} = E_t - C_{met}(T) + G_{loc}(r) + T_{coop},
\]

donde la estructura es reutilizable para cualquier ``presupuesto de estado'' local (energia, masa, momento efectivo, etc.).

Para la capa de campo, el proyecto utiliza difusion gaussiana separable con decaimiento:

\[
\phi_{t+1}(x,y) = \mathcal{D}_{3\times 3}(\phi_t)(x,y)\,(1-\lambda),\quad \lambda = 0.05.
\]

Esta forma es apropiada para representar mezclado local y disipacion de senales en la nebulosa, bajo una aproximacion de primer orden.

\section*{Marco computacional}

Implementacion actual del proyecto (verificacion directa en codigo y documentacion):

\begin{itemize}
  \item Lenguaje: Rust 2021, perfil release endurecido (LTO, \texttt{opt-level=3}, \texttt{panic=abort}).
  \item ECS: \texttt{hecs} para entidades y componentes.
  \item Paralelismo: \texttt{rayon}, buffers dobles y RNG determinista por agente/tick.
  \item Grid: 100\(\times\)100 celdas, chunking 16\(\times\)16, borde de amortiguamiento de 5 celdas.
  \item Campo: 3 canales de difusion con kernel gaussiano separable.
  \item Orquestacion: ciclo canonico de tick con actualizacion global cada 15 ticks.
  \item Metricas: exportacion JSON cada 60 ticks (Parquet por feature flag).
\end{itemize}

Presupuestos y desempeno reportados en backlog/diseno:

\begin{center}
\begin{tabular}{|l|c|c|}
\hline
Componente & Objetivo & Resultado reportado \\
\hline
Loop vacio & $<1$ ms/tick & $\sim 0.1$ ms/tick \\
Difusion 100\(\times\)100 (3 canales) & $<5$ ms/tick & $\sim 2.15$ ms/tick \\
5,000 agentes (mov+edad) & $<10$ ms/tick & $\sim 1.6$ ms/tick \\
Sistema completo (objetivo global) & $<80$ ms/tick & dentro de objetivo v1 \\
\hline
\end{tabular}
\end{center}

\section*{Consideraciones clave para construir la simulacion}

Ademas del diseno de algoritmos, la implementacion exigio un conjunto de decisiones de ingenieria que aparecen distribuidas en la documentacion de \texttt{docs/} y \texttt{design/}. Esta seccion consolida los criterios mas importantes que hubo que tener en cuenta.

\textbf{1) Alcance realista de v1 (control de complejidad).}

Desde el inicio se limito el alcance para garantizar una primera version estable y medible:

\begin{itemize}
  \item motor nativo de escritorio en lugar de despliegue web;
  \item grid 100\(\times\)100 y hasta miles de agentes, no escala masiva;
  \item difusion en CPU (sin compute shaders) para mantener trazabilidad y debug;
  \item visualizacion cientifica ligera en lugar de motor 3D inmersivo.
\end{itemize}

Esto reduce riesgo tecnico y permite validar primero la coherencia del modelo.

\textbf{2) Restricciones de hardware objetivo.}

El proyecto fue orientado a hardware de consumo (CPU 6--16 nucleos, 16--32 GB RAM). Esta condicion definio el presupuesto temporal por tick y forzo decisiones de cache-locality, paralelismo por datos y minimizacion de bloqueos.

\textbf{3) Separacion de escalas temporales.}

No todos los procesos evolucionan a la misma velocidad. Para evitar sobrecosto computacional y ruido numerico:

\begin{itemize}
  \item el loop logico corre con tick fijo;
  \item la dinamica global se actualiza cada 15 ticks;
  \item las metricas se exportan cada 60 ticks;
  \item el render opera desacoplado del hilo de simulacion.
\end{itemize}

Este desacople fue una condicion central para mantener estabilidad y rendimiento simultaneamente.

\textbf{4) Consistencia arquitectonica en capas.}

La simulacion se estructuro en capas con contratos de interfaz claros (quien lee y quien escribe cada estado). Se evito que los sistemas se llamen directamente entre si; toda comunicacion pasa por el estado global o por parametros explicitos del orquestador.

Esta regla disminuye acoplamiento, facilita pruebas y evita efectos laterales ocultos.

\textbf{5) Determinismo como requisito cientifico.}

No se considero opcional: el determinismo era necesario para reproducibilidad experimental. Para ello se adopto:

\begin{itemize}
  \item semilla global configurable por corrida;
  \item derivacion determinista por agente/tick;
  \item orden de ejecucion canonico fijo en el orquestador;
  \item uso de double-buffering en procesos de campo.
\end{itemize}

\textbf{6) Modelado de campo con estabilidad numerica.}

La difusion se implemento con kernel gaussiano separable + decaimiento y bordes absorbentes. Esta combinacion se eligio para:

\begin{itemize}
  \item evitar oscilaciones de esquemas ingenuos;
  \item controlar acumulacion artificial en bordes;
  \item mantener costo computacional bajo en CPU multicore.
\end{itemize}

\textbf{7) Diseno para observabilidad desde el inicio.}

La simulacion no se planteo como ``caja negra''. Se incluyeron metricas de emergencia, trazas y benchmarks como parte del nucleo del sistema, no como agregado tardio. Esto permitio detectar colapsos, desacoples y desviaciones spec-codigo durante el desarrollo.

\textbf{8) Gobernanza documental (spec-first).}

En \texttt{design/} se definio una regla fuerte: el comportamiento biologico debe existir primero en \texttt{simulation\_spec.md}, y la estructura de software en \texttt{architecture.md}. El codigo debe seguir estos documentos.

Este principio fue clave para coordinar trabajo por roles (Arquitecto, Biologo, Programador ECS, QA) y para reducir ambiguedad en decisiones tecnicas.

\textbf{9) Planificacion incremental por modulos.}

La hoja de ruta y backlog no priorizan cantidad de features, sino cierre de modulos con criterio de exito, benchmark y pruebas de reproducibilidad. Se avanzo de infraestructura minima a comportamiento emergente y perturbaciones en iteraciones controladas (M0--M15).

\textbf{10) Manejo explicito de trade-offs.}

Tambien hubo decisiones diferidas de forma consciente: GPU compute masivo, grids gigantes, dashboard web y evolucion genetica quedaron fuera de v1. Esta disciplina permitio proteger el objetivo principal: un simulador validable y util para experimentacion.

\section*{Algoritmos}

\textbf{1) Orquestacion por tick (orden fijo):}

\begin{enumerate}
  \item actualizacion de dinamica global cada 15 ticks;
  \item difusion de campo con swap de buffers;
  \item regeneracion de recursos;
  \item ejecucion de sistemas ECS en orden canonico;
  \item exportacion de metricas cada 60 ticks;
  \item envio no bloqueante al renderer.
\end{enumerate}

\textbf{2) Difusion numerica:}

Convolucion horizontal + vertical separables, normalizacion en bordes y decaimiento multiplicativo. Se garantiza ausencia de condiciones de carrera por lectura/escritura en buffers separados.

\textbf{3) Transiciones de estado tipo umbral:}

El proyecto implementa una regla probabilistica de tipo Fixed-Threshold:

\[
P(\text{adoptar }T)=\frac{s_T^2}{s_T^2+\theta_T^2},
\]

con heterogeneidad individual y cooldown temporal. Esta estructura es reusable para cambios de fase fisica en agregados protoplanetarios.

\textbf{4) Perturbaciones y colapso:}

Existen sistemas de perturbacion (enfermedad, depredacion en el dominio original) que modelan forzamientos adversos. En la adaptacion planetaria se interpretan como eventos disruptivos (fragmentacion, perdida de material, inestabilidad local).

\section*{Implementacion - diseno}

El diseno sigue una arquitectura en capas con contratos explicitos de lectura/escritura por sistema.

\begin{center}
\begin{tikzcd}[row sep=2em, column sep=2.2em]
\text{Orchestrator} \arrow[r] & \text{System Dynamics} \arrow[r] & \text{ECS Systems} \arrow[r] & \text{Grid/Field} \arrow[r] & \text{Metrics/Render}
\end{tikzcd}
\end{center}

Decisiones de diseno relevantes:

\begin{itemize}
  \item \textbf{Determinismo}: semilla global + RNG por agente/tick.
  \item \textbf{Rendimiento}: layout SOA para grid y chunking para localidad de cache.
  \item \textbf{Concurrencia}: sin locks en hot path; sincronizacion minima por swap.
  \item \textbf{Escalabilidad}: separacion de frecuencias (tick rapido vs dinamica lenta).
\end{itemize}

Brechas tecnicas detectadas entre especificacion y codigo actual (importantes para trazabilidad):

\begin{itemize}
  \item costo metabolico basal en codigo: \texttt{0.008/tick}; en spec: \texttt{0.02/tick};
  \item umbral de trofalaxia en codigo: \texttt{0.70}; en spec: \texttt{0.30};
  \item \texttt{initial\_honey\_reserve} por defecto en codigo: \texttt{10.0}; en spec: \texttt{0.8};
  \item fuentes de alimento: implementadas como parches fijos (17\(\times\)17), mientras la spec define fuentes puntuales/aleatorias.
\end{itemize}

Estas diferencias deben resolverse via RFC para asegurar consistencia antes de extrapolar resultados cientificos al caso planetario.

\section*{Flujo de desarrollo por iteraciones}

El desarrollo del proyecto siguio un flujo incremental guiado por \texttt{design/backlog.md}: cada iteracion cierra modulos completos con criterio de exito, pruebas y presupuesto de rendimiento. La regla operativa fue \textbf{no avanzar al siguiente bloque hasta validar el anterior}.

\textbf{Iteracion 1: Mundo y agentes minimos (M0--M3).}

Se construyo la base del motor: infraestructura Rust, loop de tick, grid chunked, difusion de feromonas y agentes ECS basicos (spawn, movimiento, edad). En esta etapa se confirmo que el nucleo corria dentro de presupuesto y con determinismo reproducible.

\textbf{Iteracion 2: Energia y mortalidad (M4--M5).}

Se incorporo el balance energetico por tick y la eliminacion de entidades sin energia. El objetivo fue cerrar el ciclo minimo de vida/muerte y garantizar limpieza correcta del estado (sin entidades ``zombie'' y sin inconsistencias de ocupacion en el grid).

\textbf{Iteracion 3: Roles y comportamiento social (M6--M9).}

Se agregaron roles diferenciados, forrajeo completo, transicion de roles por modelo de umbral fijo y trofalaxia. Esta iteracion habilito comportamiento emergente, al pasar de agentes homogeneos a una colonia funcional con especializacion y coordinacion local.

\textbf{Iteracion 4: Dinamica global y cria (M10--M11).}

Se integraron variables macro (estacionalidad, temperatura global, reservas, produccion de cria) y el ciclo huevo-larva-pupa-adulto. Con esto se acoplaron escalas micro y macro, permitiendo crecimiento y declive poblacional en funcion del estado global.

\textbf{Iteracion 5: Perturbaciones (M12--M13).}

Se implementaron amenazas internas y externas: sistema epidemiologico SIR y depredadores. Esta fase permitio evaluar resiliencia, sensibilidad a choques y condiciones de colapso bajo estres.

\textbf{Iteracion 6: Visualizacion y exportacion (M14--M15).}

Se cerraron capacidades de analisis y observacion: exportacion de metricas completas y visualizador en hilo separado. El foco fue hacer la simulacion inspeccionable en tiempo real y reutilizable en analisis posterior.

\textbf{Criterios transversales aplicados en todas las iteraciones.}

\begin{itemize}
  \item cada modulo requiere pruebas funcionales y criterio de exito explicito;
  \item se verifica rendimiento contra presupuesto antes de marcar modulo como completo;
  \item se exige reproducibilidad (misma seed, mismos resultados exportados);
  \item cualquier cambio de comportamiento debe estar documentado en spec/arquitectura.
\end{itemize}

Este flujo por iteraciones fue clave para mantener control de complejidad, trazabilidad tecnica y calidad experimental, evitando crecimiento desordenado del sistema.

\section*{Validaciones}

Validaciones implementadas en el repositorio:

\begin{itemize}
  \item pruebas unitarias de grid, difusion, movimiento, energia, mortalidad, transicion, cria, SIR y depredacion;
  \item verificacion de determinismo: mismos seeds producen mismos JSON de metricas;
  \item benchmarks con \texttt{criterion} para modulos criticos;
  \item chequeos de invariante (clamp de energia, conteo de ocupacion, bordes absorbentes).
\end{itemize}

La evidencia documental indica cobertura funcional de los modulos M0--M15 como ``completos'' en backlog tecnico.

\section*{Experimentacion (verificacion con multiples parametro)}

Se propone y documenta una matriz de experimentos compatible con la infraestructura actual:

\begin{center}
\begin{tabular}{|c|c|c|c|c|}
\hline
Escenario & Semilla & Perturbacion & Poblacion inicial & Objetivo \\
\hline
E0 (base) & 42 & sin perturbacion extra & 500 & estabilidad y baseline \\
E1 (disruptivo) & 44 & enfermedad habilitada & 500 & robustez a contagio \\
E2 (choques externos) & 99 & depredadores $>0$ & 500 & resiliencia a perdidas \\
E3 (estres extremo) & 42 & baja reserva inicial & 10--100 & tiempo a colapso \\
\hline
\end{tabular}
\end{center}

Indicadores a medir en cada corrida:

\begin{itemize}
  \item \texttt{population\_by\_role},
  \item \texttt{colony\_reserve},
  \item \texttt{pheromone\_entropy},
  \item \texttt{sir\_prevalence},
  \item \texttt{time\_to\_collapse},
  \item \texttt{mortality\_rate} desagregada.
\end{itemize}

En una corrida observada de larga duracion (archivos de metricas existentes):

\begin{itemize}
  \item en tick 0 se observa poblacion total 500 y produccion activa;
  \item se registra colapso temprano en \texttt{time\_to\_collapse = 86};
  \item para ticks tardios (ej. 1800 y 3720) no quedan agentes activos, pero la reserva global permanece positiva.
\end{itemize}

\section*{Analisis de resultados con base en experimentacion}

Interpretacion tecnica:

\begin{itemize}
  \item el motor es computacionalmente estable y reproducible, por lo que es apto para barridos parametrizados;
  \item la dinamica emergente muestra sensibilidad alta a parametros de consumo, reposicion y perturbacion;
  \item la existencia de colapso temprano junto con reserva positiva sugiere desacople entre subsistemas locales y macroestado;
  \item las brechas spec-codigo afectan la interpretacion causal de resultados y deben cerrarse antes de validar hipotesis cientificas fuertes.
\end{itemize}

En el contexto ``formacion de planetas'', este hallazgo se traduce en una advertencia metodologica: el simulador puede producir trayectorias de colapso o estabilizacion que dependen mas de la parametrizacion interna que del fenomeno fisico objetivo si no se mantiene consistencia formal entre especificacion e implementacion.

\section*{Conclusiones}

\begin{itemize}
  \item BeeFlow entrega una base solida para simulacion hibrida multi-escala: ECS + grid + dinamica global + metricas.
  \item El proyecto ya cumple requisitos clave de ingenieria cientifica: determinismo, observabilidad, rendimiento y modularidad.
  \item Para usarlo como plataforma de estudio de formacion planetaria, el siguiente paso critico es una capa de parametrizacion fisica explicita y una alineacion estricta spec-codigo.
  \item La experimentacion inicial muestra que el framework detecta regmenes de estabilidad/colapso y es adecuado para analisis de sensibilidad con multiples parametros.
\end{itemize}

\textbf{Recomendacion inmediata:} abrir RFC de armonizacion de parametros nucleares (costos, umbrales, condiciones de fuente) y luego ejecutar una campana de experimentos factorial con al menos 30 semillas por escenario.

\end{document}
